\begin{explanationbox}
  \textbf{\textcolor{CExplanation}{先想一个问题}}
  原来 $y=x^2$，把图像上每个点的横坐标\textbf{缩小一半}，点 $(2,4)$ 就移到 $(1,4)$。那么新图像经过 $(1,4)$，新式子该怎样写？

  \medskip
  \textbf{\textcolor{CExplanation}{用同一个高度倒着找横坐标}}
  在原函数中，要得到高度 $4$，可取 $x=2$；在新函数中，同样的高度 $4$ 出现在 $x=1$。\textbf{新横坐标只有原来的一半，代入原函数时就需要把它乘 $2$：}
  \[
    \underbrace{f(2\times 1)}_{\text{新图像在 }x=1\text{ 的高度}}
    =\underbrace{f(2)}_{\text{原图像在 }x=2\text{ 的高度}}=4.
  \]
  因此 $y=f(x)$ 压窄一半后变成 $y=f(2x)$，对于 $f(x)=x^2$ 就是 $y=4x^2$。

  \begin{center}
  \begin{tikzpicture}[x=1.12cm,y=0.66cm,>=stealth,every node/.style={font=\small}]
    \draw[->,gray!75] (-0.15,0)--(4.65,0) node[right] {横坐标};
    \foreach \x in {0,1,2,3,4}{\draw (\x,0.06)--(\x,-0.06) node[below] {$\x$};}
    \draw[blue!75!black,very thick] (0,1.1)--(4,1.1);
    \fill[blue!75!black] (2,1.1) circle (2.3pt);
    \fill[orange!85!black] (1,1.1) circle (2.3pt);
    \draw[->,thick,orange!85!black] (1.89,1.35)--(1.10,1.35);
    \node[above,blue!75!black] at (2.45,1.1) {原点 $(2,4)$};
    \node[above,orange!85!black] at (0.95,1.55) {新点 $(1,4)$};
    \node[right] at (4.08,1.1) {高度同为 $4$};
  \end{tikzpicture}

  \end{center}

  \textbf{\textcolor{CExplanation}{推广：不用死记，用一个点来检查}}
  若原图像上有点 $(a,f(a))$，横坐标变为原来的 $k$ 倍后，对应点是 $(ka,f(a))$。要让新函数在 $x=ka$ 时仍算出 $f(a)$，就应使用
  \[
     g(x)=f\!\left(\frac{x}{k}\right),\qquad
     g(ka)=f(a).
  \]
  这就是\textbf{“横坐标乘 $k$，括号里除以 $k$”}的原因；写作 $f(\omega x)$ 时，实际横向倍数就是 $1/\omega$。

  \medskip
  \textbf{\textcolor{CExplanation}{考试什么时候最快用？}}
  \begin{itemize}[leftmargin=*,itemsep=0.3em,topsep=0.3em]
    \item \textbf{求新解析式：} 题目说“横坐标变成 $3$ 倍”，直接把原函数中的 $x$ 换成 $x/3$。
    \item \textbf{判断图像宽窄：} 看括号里的 $x$：乘 $3$ 就横向压成 $1/3$，乘 $1/2$ 就拉成 $2$ 倍。
    \item \textbf{检查答案：} 任选原图像上的一个点，看新坐标的\textbf{高度有没有变化}。
  \end{itemize}
  \textbf{补充提醒：} 这种变换只改变横坐标。若原函数的定义域为 $D$，新函数 $f(x/k)$ 的定义域为所有满足 $x/k\in D$ 的 $x$；\textbf{不是所有函数都能随意取负数或零作输入}。
\begin{center}
\begin{tikzpicture}[font=\small,>=stealth]
\draw[draw=teal!50!black,fill=green!3,rounded corners=3pt] (0,0) rectangle (13,7.2);
\node at (6.5,6.6) {压缩不是整体左移：负横坐标点向右，正横坐标点向左};
\begin{scope}[shift={(3.3,1.6)},x=.85cm,y=.55cm]
\draw[gray,->] (-3.3,0)--(3.5,0) node[right] {$x$};\draw[gray,->] (0,-.3)--(0,7.3) node[above] {$y$};
\draw[teal!70!black,thick,domain=-3:3.2] plot(\x,{\x+3});
\draw[orange!85!black,thick,domain=-1.5:1.6] plot(\x,{2*\x+3});
\foreach \old/\new/\height in {-2/-1/1,2/1/5}{
 \draw[->,blue!65!black,very thick] (\old,\height)--(\new,\height);
 \fill[teal!70!black] (\old,\height) circle(2.5pt);\fill[orange!85!black] (\new,\height) circle(2.5pt);
}
\node[below] at (-2,.85) {$A$};\node[below right] at (-1,.85) {$A'$};
\node[below right] at (2,5) {$B$};\node[above left] at (1,5) {$B'$};
\end{scope}
\node[anchor=west,align=left,text width=4.5cm] at (8.1,3.6) {绿线 $f(x)=x+3$。\\[8pt]橙线 $f(2x)=2x+3$。\\[8pt]$A(-2,1)\to A'(-1,1)$\\[8pt]$B(2,5)\to B'(1,5)$\\[8pt]横坐标减半，高度不变。};
\node at (6.5,.45) {看左右两侧的对应点，才能分清“压窄”和“平移”。};
\end{tikzpicture}

\end{center}
\textbf{两侧一起看：}压缩时左右两侧都靠近 $y$ 轴；这与所有点向同一方向平移不同。
\end{explanationbox}
